% Figure: a second-order band-pass transfer function read off the
% complex plane. Left panel: the pole-zero plot in the s-plane (a zero at
% the origin, a conjugate pair of poles in the left half-plane). Right
% panel: the resulting magnitude response, the distance ratio read off
% as the test point i*omega slides up the imaginary axis.
\begin{figure}[htbp]
\centering
\begin{tikzpicture}[line cap=round, >=latex, font=\footnotesize]
  % ---------- Left: pole-zero plot ----------
  \begin{scope}[scale=1.0]
    \draw[->, engblack, thin] (-2.6,0) -- (0.7,0) node[right] {$\sigma$};
    \draw[->, engblack, thin] (0,-2.4) -- (0,2.5) node[above] {$j\omega$};
    % pole pair at sigma = -0.4, omega = +-1.8
    \node[engred] at (-0.4,1.8) {$\times$};
    \node[engred] at (-0.4,-1.8) {$\times$};
    \node[engred, anchor=west] at (-0.32,2.05) {pole};
    % zero at origin
    \node[engblue] at (0,0) {$\circ$};
    \node[engblue, anchor=north west] at (0.05,-0.05) {zero};
    % test point on imaginary axis
    \filldraw[engblack] (0,1.5) circle (0.04);
    \node[engblack, anchor=west] at (0.07,1.5) {$j\omega$};
    % distance vectors
    \draw[engred!70, thin, ->] (-0.4,1.8) -- (0,1.5);
    \draw[engblue!70, thin, ->] (0,0) -- (0,1.5);
    \node[anchor=south] at (-1.0,-2.3) {$s$-plane};
  \end{scope}
  % ---------- Right: magnitude response ----------
  \begin{scope}[shift={(5.2,-2.2)}, scale=1.0]
    \draw[->, engblack, thin] (0,0) -- (4.4,0) node[right] {$\omega$};
    \draw[->, engblack, thin] (0,0) -- (0,3.2) node[above] {$|H|$};
    % resonant peak near omega0 = 1.8 (x = 1.8 in these units, scaled)
    \draw[engblue, thick] plot[domain=0.05:4.2, samples=120, smooth]
      ({\x}, {2.6/sqrt( (1 + ( (\x*\x - 3.24)/(0.8*\x) )^2 ) )});
    \draw[engblack!40, dashed] (1.8,0) -- (1.8,2.6);
    \node[anchor=north] at (1.8,-0.05) {$\omega_0$};
  \end{scope}
\end{tikzpicture}
\caption[Reading a transfer function off the complex plane]{A
second-order band-pass response, read directly off the pole-zero plot.
The transfer function $H(s) = K s / [(s - p)(s - \bar p)]$ has a zero at
the origin and a conjugate pole pair at $p = -\sigma_0 + j\omega_0$. The
magnitude at frequency $\omega$ is the product of distances from the test
point $j\omega$ to each zero divided by the product of distances to each
pole. As $j\omega$ slides up the imaginary axis past the pole's
imaginary part $\omega_0$, the near-pole distance shrinks to its minimum
and the magnitude peaks: the resonance is the geometric event of the
test point passing closest to a pole. The damping $\sigma_0$ (the pole's
distance from the imaginary axis) sets the sharpness of the peak.}
\label{fig:vol02:ch03:pole-zero-bode}
\end{figure}
